% =============================================================================
\section{Governing Equations}
\label{sec:dae}
% =============================================================================

% -----------------------------------------------------------------------------
\subsection{System State Vector}
\label{sec:dae:state}
% -----------------------------------------------------------------------------

Each tank node $i \in \{1,\dots,N\}$ carries three dynamic states:
\begin{equation}
  \vect{x}_i = \begin{bmatrix} m_i \\ T_i \\ T_{s,i} \end{bmatrix},
\end{equation}
where $m_i$ is the hydrogen mass [\si{\kilogram}], $T_i$ is the bulk gas
temperature [\si{\kelvin}], and $T_{s,i}$ is the lumped solid (wall/liner)
temperature [\si{\kelvin}].

The full system state vector is
\begin{equation}
  \vect{x} = \begin{bmatrix}
    \vect{x}_1 \\ \vdots \\ \vect{x}_N
  \end{bmatrix} \in \mathbb{R}^{3N}.
\end{equation}

The specific volume and density are derived quantities:
\begin{equation}
  V_i = \text{const.},\qquad
  \rho_i = \frac{m_i}{V_i}.
\end{equation}
All thermodynamic properties ($P_i$, $h_i$, $c_{v,i}$, \dots) are evaluated
from $(T_i, \rho_i)$ via CoolProp.

% -----------------------------------------------------------------------------
\subsection{Conservation Equations}
\label{sec:dae:ode}
% -----------------------------------------------------------------------------

\subsubsection{Mass balance}

Let $\Edges_i^+$ and $\Edges_i^-$ denote the sets of edges entering and leaving
node $i$, respectively.  The mass balance is
\begin{equation}
  \dot{m}_i = \sum_{e \in \Edges_i^+} \mdot_e - \sum_{e \in \Edges_i^-} \mdot_e
            \equiv \mdot_i^{\text{net}}.
  \label{eq:mass-balance}
\end{equation}

\subsubsection{Gas temperature}

For an isochoric open system, the first law of thermodynamics gives (see
derivation in \cref{app:energy})
\begin{equation}
  m_i c_{v,i}\, \dot{T}_i =
    \underbrace{\sum_{e \in \Edges_i^+} \mdot_e \bigl(h_e - h_i\bigr)
      - \sum_{e \in \Edges_i^-} \mdot_e \bigl(h_e - h_i\bigr)}_{\text{net enthalpy flux}}
    + \underbrace{\frac{T_i}{\rho_i}\,\alpha_i\,\mdot_i^{\text{net}}}_{\text{real-gas work term}}
    + \underbrace{\dot{Q}_{s,i}}_{\text{wall heat}},
  \label{eq:energy-balance}
\end{equation}
where
\begin{align}
  \alpha_i &= \at{\frac{\partial P}{\partial T}}{\rho}
    \quad \text{(isochoric pressure--temperature coefficient)}, \\
  h_e      &= \text{enthalpy of the fluid carried by edge } e
    \quad \text{(evaluated at the source node)}, \\
  h_i      &= h(T_i, \rho_i)
    \quad \text{(local specific enthalpy)}.
\end{align}
For discharge and vent edges leaving tank $i$, the fluid enthalpy equals the
local value ($h_e = h_i$), so those terms vanish from the enthalpy flux.

\subsubsection{Solid temperature}

The liner and composite wall are modelled as a single lumped thermal mass:
\begin{equation}
  \dot{T}_{s,i} = \frac{\dot{Q}_{s,i}^{\text{ext}} - \dot{Q}_{s,i}}{m_{s,i}\,c_{s,i}},
  \label{eq:solid-temp}
\end{equation}
where $\dot{Q}_{s,i}^{\text{ext}}$ is the external (ambient) heat influx computed
by the thermal model, and $\dot{Q}_{s,i}$ is the gas-side heat transfer.  The
thermal model follows the Stops convection formulation.

% -----------------------------------------------------------------------------
\subsection{Edge Flow-Rate Models}
\label{sec:dae:edges}
% -----------------------------------------------------------------------------

\subsubsection{Coupling edge}

Coupling valves are modelled as one of several valve archetypes
(\texttt{PressureTriggeredValve}, \texttt{FeedforwardPressureEnforcer}, etc.).
The generic interface is
\begin{equation}
  \mdot_e^{\text{coup}} = f_\text{valve}(P_\text{src},\, P_\text{tgt},\, t;\; \theta),
\end{equation}
where $\theta$ denotes valve-specific parameters (orifice diameter, opening
pressure, PID gains, \dots).  The enthalpy of the transferred mass is evaluated
at the source node: $h_e = h(T_\text{src}, \rho_\text{src})$.  If a peripheral
component chain is declared on the edge (\cref{sec:network:peripheral}), the
chain transforms $h_e$ before it enters the downstream energy balance.

\subsubsection{Discharge edge}

In the normal operating regime (Configuration~A, \cref{sec:dae:config}), the
discharge flow rate is prescribed by the mission profile:
\begin{equation}
  \mdot_e^{\text{disc}} = \mdot_\text{CSV}(t) \quad [\text{Config~A}].
\end{equation}
In Configuration~B the discharge rate continues at the mission value;
a compensating heat flux $\dot{Q}_{HX,i}$ is supplied to maintain the
pressure floor (Section~\ref{sec:dae:configB}).

\subsubsection{Vent edge}

Venting is activated when the tank pressure reaches the maximum operating
pressure $P_{\text{vent},i}$.  The flow rate is
\begin{equation}
  \mdot_e^{\text{vent}} = k_v\,\max\!\bigl(0,\, P_i - P_{\text{vent},i}\bigr),
  \label{eq:vent-rate}
\end{equation}
with proportionality constant $k_v = \SI{1e-8}{\kilogram\per\second\per\pascal}$.
Only the flow rate is needed for venting (per \cref{sec:dae:config}, Config~C).

\subsubsection{Refuel edge}

For refuel missions the incoming mass is compressed from a low-pressure source
reservoir by an ideal cryogenic pump with isentropic efficiency $\eta_p = 0.78$:
\begin{equation}
  h_e^{\text{cryo}}(P_i) = h_1 + \frac{h_{2s} - h_1}{\eta_p},\quad
  h_1 = h(P_0, x{=}0),\quad
  h_{2s} = h(P_i, s_1),
\end{equation}
where $P_0 = \SI{3}{\bar}$ is the suction pressure and $s_1$ is the specific
entropy at suction conditions.

% -----------------------------------------------------------------------------
\subsection{Operating Configurations}
\label{sec:dae:config}
% -----------------------------------------------------------------------------

Each tank independently operates in one of three configurations depending on its
current pressure:

\begin{center}
\begin{tabularx}{\linewidth}{@{} c l X @{}}
\toprule
\textbf{Config} & \textbf{Activation condition} & \textbf{Effect} \\
\midrule
A & $P_{\min,i} < P_i < P_{\text{vent},i}$ &
  Free ODE evolution.  All edge flows as prescribed. \\[4pt]
B & $P_i \le P_{\min,i}$ &
  Pressure-floor constraint: $\dot{P}_i = 0$.
  Discharge continues at the prescribed mission rate; a compensating heat
  flux $\dot{Q}_{HX,i}$ is supplied by a heat exchanger to maintain
  constant pressure
  (Section~\ref{sec:dae:configB}). \\[4pt]
C & $P_i \ge P_{\text{vent},i}$ &
  Vent edge activated; $\mdot_e^{\text{vent}} > 0$. \\
\bottomrule
\end{tabularx}
\end{center}

Config~B is implemented with a hysteresis band to prevent chattering.  The
transition $\text{A}\to\text{B}$ occurs at $P_i \le P_{\min,i}(1-\delta)$ and
the return $\text{B}\to\text{A}$ at $P_i > P_{\min,i}(1+\delta)$, where
$\delta = 0.01$.

% -----------------------------------------------------------------------------
\subsection{Configuration~B: Compensating Heat Flux}
\label{sec:dae:configB}
% -----------------------------------------------------------------------------

When $P_i = P_{\min,i}$, pressure must not fall further while discharge
continues at the mission-prescribed rate $\mdot_\text{CSV}(t)$.  A compensating
heat flux $\dot{Q}_{HX,i}$ is supplied to the gas by a heat exchanger coupled
to the tank (for example, waste heat from the fuel cell, ambient heating, or a
dedicated heater).  No discharge throttling occurs.

Imposing $\dot{P}_i = 0$ on $P_i = P(T_i, \rho_i)$:
\begin{equation}
  \alpha_i\,\dot{T}_i + \frac{\beta_i}{V_i}\,\mdot_i^{\text{net}} = 0,
  \label{eq:config-b-constraint}
\end{equation}
where $\beta_i = \at{\partial P/\partial\rho}{T}$.  Solving for $\dot{T}_i$
and substituting into the energy balance \cref{eq:energy-balance}:
\begin{equation}
  m_i c_{v,i}
    \left(-\frac{\beta_i}{\alpha_i}\frac{\mdot_i^{\text{net}}}{V_i}\right)
  = \sum_{e\in\Edges_i^+}\!\mdot_e(h_e - h_i)
    + \frac{T_i}{\rho_i}\alpha_i\,\mdot_i^{\text{net}}
    + \dot{Q}_{s,i} + \dot{Q}_{HX,i}.
  \label{eq:config-b-energy}
\end{equation}
Applying the triple-product identity
$\at{\partial T/\partial\rho}{P} = -\beta_i/\alpha_i$ and rearranging for
$\dot{Q}_{HX,i}$ (pure-discharge case, $\Edges_i^+ = \varnothing$):
\begin{equation}
  \boxed{
  \dot{Q}_{HX,i} = \mdot_{\text{out},i}
    \underbrace{\!\left(
      \frac{T_i}{\rho_i}\,\alpha_i
      - \rho_i\,c_{v,i}\at{\frac{\partial T}{\partial\rho}}{P}
    \right)}_{>0}
    - \dot{Q}_{s,i}
  }
  \label{eq:config-b-qhx}
\end{equation}
Both derivatives are evaluated from CoolProp at $(T_i, \rho_i)$:
$\alpha_i = \texttt{d(P)/d(T)|D}$ and
$\at{\partial T/\partial\rho}{P} = \texttt{d(T)/d(D)|P}$.
When inflows are present, the right-hand side acquires an additional
$-\sum_{e\in\Edges_i^+}\mdot_e(h_e - h_i)$ term.

\paragraph{Physical interpretation.}
\Cref{eq:config-b-qhx} states that the heat exchanger must compensate two
competing effects: (i)~the thermodynamic pressure--temperature coupling
$(T/\rho)\alpha_i$ that would cause pressure to fall as gas leaves the tank, and
(ii)~the isobaric cooling associated with the density change,
$\rho_i c_{v,i}\at{\partial T/\partial\rho}{P}$.  Because
$\at{\partial T/\partial\rho}{P} < 0$ for hydrogen under typical cryogenic
and high-pressure conditions, both terms are positive and $\dot{Q}_{HX,i} > 0$
(heat is \emph{added}, not removed).  In the limit of negligible ambient heat
leak ($\dot{Q}_{s,i} \approx 0$) the required heat input is proportional to the
discharge rate.

